% !PRINTED 374
\weberpage{374}
% Source-faithful Cayley configuration signs: a four-stroke star and a seven-rayed star.
\providecommand{\weberfourup}{\mathord{\tikz[baseline=-0.2ex,x=0.23em,y=0.23em,line width=0.45pt]{\draw (0,0)--(-1.05,1) (0,0)--(-0.35,1.08) (0,0)--(0.35,1.08) (0,0)--(1.05,1);}}}
\providecommand{\webersevenstar}{\mathord{\tikz[baseline=-0.35ex,x=0.21em,y=0.21em,line width=0.45pt]{\draw (0,0)--(90:1) (0,0)--(38.57:1) (0,0)--(-12.86:1) (0,0)--(-64.29:1) (0,0)--(-115.71:1) (0,0)--(-167.14:1) (0,0)--(141.43:1);}}}
%<<CONT
dieser Bestimmungen ist daher 280, und wir kommen zu folgendem Resultate:

\gesperrt{Es giebt im Ganzen 288 Aronhold'sche Systeme, deren Zeichen eine der beiden Gestalten hat}

\[
\webersevenstar;\ \Delta\,\weberfourup.
\]

Aus diesen Zeichen darf man aber nicht etwa schliessen, dass diese Siebener-Systeme in zwei verschiedene Arten zerfallen. Der Unterschied der beiden Figuren liegt nur in der Bezeichnung. In der That können wir ja von einem ganz beliebigen der vollständigen Siebener-Systeme ausgehen, um die Bezeichnung abzuleiten.

Hiernach findet man leicht die Bezeichnung für die \gesperrt{Steiner'schen} Complexe, die nach einem der beiden folgenden Typen zu bilden sind:

\[
[1\,2][3\,4],\ [1\,3][2\,4],\ [1\,4][2\,3],\ [5\,6][7\,8],\ [5\,7][6\,8],\ [5\,8][6\,7],
\]
\[
[1\,7][1\,8],\ [2\,7][2\,8],\ [3\,7][3\,8],\ [4\,7][4\,8],\ [5\,7][5\,8],\ [6\,7][6\,8].
\]

\begin{center}
\S.\,100.

\gesperrt{Rationale Bestimmung der Curve aus einem vollständigen Siebener-Systeme.}
\end{center}

Die grosse Bedeutung der \gesperrt{Aronhold'schen} Systeme für das Problem der Doppeltangenten spricht sich in folgenden beiden Sätzen aus:

I. \gesperrt{Ist bei einer Curve vierter Ordnung ein vollständiges Siebener-System gegeben, so können daraus alle übrigen Doppeltangenten rational bestimmt werden.}

II. \gesperrt{Sind sieben beliebige gerade Linien in einer Ebene gegeben, so kann man im Allgemeinen, d.\,h. wenn gewisse rationale Functionen von den Co\"efficienten in den Gleichungen dieser Geraden nicht verschwinden, auf rationalem Wege die Gleichung einer Curve vierter Ordnung ohne singulären Punkt bestimmen, für die die gegebenen sieben Linien ein Aronhold'sches System bilden.}

Um den ersten Satz zu beweisen, würde es bei der vollkommenen Gleichberechtigung der Ziffern 1 bis 7 genügen, die
%CONT>>
